\chapter{Rezolvarea recurențelor liniare prin exponențiere rapidă}

Să pornim de la codul binecunoscut pentru calculul lui $a^n$ în $\bigoh(\log n)$:

\begin{minted}{c}
int mod_pow(int a, int n) {
  long long result = 1;

  while (n) {
    if (n & 1) {
      result = result * a % MOD;
    }
    a = (long long)a * a % MOD;
    n >>= 1;
  }

  return result;
}
\end{minted}

Algoritmul funcționează rapid pentru că accelerează succesiuni lungi de înmulțiri simple, înlocuindu-le cu ridicări la pătrat. Putem folosi aceeași rețetă pentru a ridica la putere orice obiecte care pot fi înmulțite și a căror înmulțire este asociativă, în particular matricele pătrate și permutările. Trebuie doar ca, în loc de o valoare întreagă, \ccode{a} să fie o matrice, iar în loc de operatorul \ccode{*} pentru înmulțire de întregi să implementăm înmulțirea de matrice.

Vom vedea cum exponențierea de matrice ne poate ajuta în probleme de programare dinamică în care trebuie să împingem limita pentru $n$ pînă la valori foarte mari.

\section{Șirul lui Fibonacci. Matrice generatoare}

Să considerăm un exemplu celebru: șirul lui Fibonacci, definit ca

$$
\begin{cases}
  F_0 &= 0 \\
  F_1 &= 1 \\
  F_n &= F_{n-1} + F_{n-2} \quad \textrm{pentru} \; n \geq 2
\end{cases}
$$

Dacă $n$ este foarte mare, de ordinul a $10^{18}$, atunci calculul liniar nu ne mai ajută. Vom încerca să exprimăm recurența cu următoarele elemente:

\begin{itemize}
  \item Un vector-coloană, $V_n$, pentru pasul curent al formulei. Acesta trebuie să includă cel puțin termenul de calculat, $F_n$.
  \item Un vector-coloană, $V_{n-1}$, pentru pasul anterior al formulei. Acesta trebuie să includă toți termenii de care depine $F_n$, în cazul nostru $F_{n-1}$ și $F_{n-2}$, și posibil și alte cantități.
  \item O matrice $A$, numită \textbf{matrice generatoare}, cu proprietatea că
\end{itemize}

$$V_{n} = A \cdot V_{n-1}$$

Pentru șirul lui Fibonaci, obținem relația:

$$
\begin{pmatrix} F_{n} \\ F_{n-1} \end{pmatrix}
=
\begin{pmatrix} 1 & 1 \\ 1 & 0 \end{pmatrix}
\begin{pmatrix} F_{n-1} \\ F_{n-2} \end{pmatrix}
$$

Remarcăm că matricea este constantă (independentă de $n$). De aceea, putem scrie:

$$
\begin{pmatrix} F_{n} \\ F_{n-1} \end{pmatrix}
=
\begin{pmatrix} 1 & 1 \\ 1 & 0 \end{pmatrix}^{n-1}
\begin{pmatrix} F_{1} \\ F_{0} \end{pmatrix}
$$

Iar pe $A^{n-1}$ îl putem calcula cu exponențiere rapidă.

Arta pentru acest gen de probleme este tocmai să găsim forma vectorului-coloană și matricea generatoare. Observăm că vectorul-coloană din dreapta trebuie să conțină toate cantitățile necesare pentru calculul lui $F_n$, ceea ce dictează și structura vectorului din stînga (cei doi sînt identici, doar indicii diferă cu 1).

Să vedem două exemple mai complicate.

\section{Alte recurențe liniare}

Să considerăm un exemplu în care recurența implică 4 termeni anteriori. Vom folosi o matrice de $4 \times 4$. Termenii $F_0$, $F_1$, $F_2$ și $F_3$ sînt dați, iar formula de recurență este:

$$F_{n} = a F_{n - 1} + b F_{n - 2} + c F_{n - 3} + d F_{n-4}$$

Atunci recurența exprimată matricial este:


$$
\begin{pmatrix} F_{n} \\ F_{n-1} \\ F_{n-2} \\ F_{n-3} \end{pmatrix}
=
\begin{pmatrix}
  a & b & c & d \\
  1 & 0 & 0 & 0 \\
  0 & 1 & 0 & 0 \\
  0 & 0 & 1 & 0 \\
\end{pmatrix}
\begin{pmatrix} F_{n-1} \\ F_{n-2} \\ F_{n-3} \\ F_{n-4} \end{pmatrix}
$$

\section{Recurențe liniare cu termen constant}

Să considerăm acum recurența:

$$F_{n} = a F_{n - 1} + b F_{n - 2} + c F_{n-3} + d$$

Secretul aici este să nu încercăm să construim vectori-coloană doar de 3 elemente. Dat fiind că există și un termen liber, trebuie să punem și o valoare scalară pe vectorii-coloană. Avem două variante:

$$
\begin{pmatrix} F_{n} \\ F_{n-1} \\ F_{n-2} \\ 1 \end{pmatrix}
=
\begin{pmatrix}
  a & b & c & d \\
  1 & 0 & 0 & 0 \\
  0 & 1 & 0 & 0 \\
  0 & 0 & 0 & 1 \\
\end{pmatrix}
\begin{pmatrix} F_{n-1} \\ F_{n-2} \\ F_{n-3} \\ 1 \end{pmatrix}
$$

sau

$$
\begin{pmatrix} F_{n} \\ F_{n-1} \\ F_{n-2} \\ d \end{pmatrix}
=
\begin{pmatrix}
  a & b & c & 1 \\
  1 & 0 & 0 & 0 \\
  0 & 1 & 0 & 0 \\
  0 & 0 & 0 & 1 \\
\end{pmatrix}
\begin{pmatrix} F_{n-1} \\ F_{n-2} \\ F_{n-3} \\ d \end{pmatrix}
$$
